\documentclass[12pt]{article}
\usepackage[utf8]{inputenc}
\usepackage[T1]{fontenc}
\usepackage[margin=1in]{geometry}
\usepackage{setspace}
\usepackage{xcolor}
\usepackage[colorlinks=true, urlcolor=blue!60!black, linkcolor=blue!60!black]{hyperref}
\setlength{\parindent}{2em}
\setlength{\parskip}{1.25em}
\renewcommand{\baselinestretch}{1.0}
\begin{document}
\begin{center}
{\LARGE \textit{Lujan v. Defenders of Wildlife}, 504 U.S. 555 (1992)\par}
\vspace{0.8em}
{\large Case Brief --- Four-Page Edition\par}
\end{center}
\vspace{1em}

\section{Memory Jogger}

Establishing the modern three-part Article III standing test, the Court held that environmental plaintiffs with vague ``some day'' plans to revisit foreign habitats lacked ``actual or imminent'' injury, and that Congress cannot confer standing through a citizen-suit provision to vindicate a generalized procedural grievance (\href{https://supreme.justia.com/cases/federal/us/504/555/}{Justia}).

\section{Detailed Case Facts}

Section 7(a)(2) of the Endangered Species Act requires each federal agency, in consultation with the Secretary of the Interior, to ensure that any action it authorizes, funds, or carries out is not likely to jeopardize endangered species. A 1978 regulation applied that duty to federal actions abroad; a 1986 regulation limited it to the United States and the high seas.

Defenders of Wildlife sued the Secretary under the ESA's citizen-suit provision to restore the 1978 interpretation. For standing it relied on two members' affidavits. Joyce Kelly had visited Egypt in 1986, observed the habitat of the endangered Nile crocodile, and intended to return; she objected to U.S. involvement in rehabilitating the Aswan High Dam. Amy Skilbred had visited Sri Lanka in 1981 near the AID-funded Mahaweli project without seeing the endangered elephants or leopards, and intended to return ``in the future,'' but admitted she had no current plans because of Sri Lanka's civil war.

\section{Procedural History}

The district court dismissed for lack of standing; a divided Eighth Circuit reversed. On remand the district court granted summary judgment to Defenders on the merits, and the Eighth Circuit affirmed, finding standing (including a ``procedural injury'') and holding that \S~7(a)(2) applies abroad. The Supreme Court reversed on June 12, 1992.

\section{Judicial Votes}

Reversed, \textbf{7--2} in the judgment; \textbf{6--3} that respondents lacked standing.
\begin{itemize}
\item \textbf{Opinion of the Court (Parts I, II, III-A, IV):} Justice Scalia, joined by Chief Justice Rehnquist and Justices White, Kennedy, Souter, and Thomas.
\item \textbf{Plurality (Part III-B, redressability):} Scalia, Rehnquist, White, and Thomas.
\item \textbf{Concurring in part and in the judgment:} Justice Kennedy, joined by Justice Souter.
\item \textbf{Concurring in the judgment:} Justice Stevens (standing existed, but \S~7(a)(2) does not apply abroad).
\item \textbf{Dissenting:} Justice Blackmun, joined by Justice O'Connor.
\end{itemize}

\section{Holding}

The ``irreducible constitutional minimum'' of standing requires (1) an injury in fact that is concrete, particularized, and actual or imminent; (2) causation fairly traceable to the defendant rather than third parties; and (3) likely redressability by a favorable decision, each proved with the evidence required at that stage of litigation. Congress cannot confer standing by authorizing ``any person'' to enforce an abstract procedural duty unconnected to a concrete injury. Part III-B's redressability analysis is only a plurality view, and Footnote 7's relaxed rule for procedural rights tied to concrete interests is influential dicta.

\section{Analysis of Opinions}

\textbf{Majority and plurality (Scalia).} Standing derives from the separation of powers and is ``substantially more difficult'' to establish when the plaintiff is not the object of the regulation. Wildlife observation is a cognizable interest, but past visits prove nothing, and ``some day'' intentions to return fail. The ``ecosystem,'' ``animal,'' and ``vocational'' nexus theories would stretch standing ``into pure speculation and fantasy.'' The plurality added that the funding agencies were not parties and supplied only a fraction of the projects' funding (AID less than 10\% of Mahaweli). Part IV held that letting Congress convert the public interest in lawful administration into an individual right would transfer the President's duty to ``take Care that the Laws be faithfully executed'' to the courts.

\textbf{Kennedy (with Souter).} He would not foreclose nexus theories on a different record and emphasized that Congress may define new injuries, so long as it identifies the injury and relates it to the class of persons entitled to sue.

\textbf{Stevens.} Imminence should turn on the timing of the environmental harm, not the plaintiff's travel plans, and agencies would follow the Court's ruling; but he concurred because \S~7(a)(2) does not apply abroad.

\textbf{Blackmun (with O'Connor).} The majority weighed evidence at summary judgment and demanded plane tickets, ``an empty formality.'' He accused the Government of playing ``Three-Card Monte'' on redressability and called Part IV ``an anachronistically formal view of the separation of powers.''

\section{Examples: Future Applications}

\begin{enumerate}
\item \textbf{Same-Side (No Standing): Hopeful Photographer.} A photographer who visited a Brazilian rainforest five years ago and hopes to return ``someday'' cannot challenge a federal highway loan.
\item \textbf{Same-Side (No Standing): Any-Citizen Deadline Suit.} A citizen with no personal use for an agency report cannot sue under an ``any U.S. citizen'' provision merely to compel timely publication.
\item \textbf{Opposite-Side (Standing): Downstream Homeowner.} A homeowner below a federally licensed dam may challenge the missing environmental impact statement, under Footnote 7.
\item \textbf{Opposite-Side (Standing): Scheduled Researcher.} A marine biologist with a grant, reservations, and permits for next month's study of the very reef to be dredged shows imminent, site-specific injury.
\item \textbf{Fence-Sitter: Undisseminated Credit Error.} A consumer with an inaccurate but never-shared credit alert invokes a statutory right; Kennedy's view suggests standing, Scalia's concrete-injury demand suggests not---the tension later addressed in \textit{Spokeo} and \textit{TransUnion}.
\end{enumerate}

\section{Critique}

\textbf{Originalist views.} Defenders, building on Scalia's 1983 article on standing and the separation of powers, praise the decision for confining judges to private-rights adjudication. But historical scholars note that ``injury in fact'' is a twentieth-century invention and that early Congresses enacted informer and \textit{qui tam} statutes; Justice Thomas has urged a private-rights/public-rights framework instead.

\textbf{Progressive views.} Cass Sunstein argued that invalidating a congressionally authorized citizen suit lacked support in Article III's text or history; Richard Pierce called the ``plane ticket'' rule formalistic and manipulable; environmental scholars argue the test disadvantages diffuse ecological harms.

\section{Key Quotations}

\begin{quote}
``Over the years, our cases have established that the irreducible constitutional minimum of standing contains three elements. First, the plaintiff must have suffered an `injury in fact'---an invasion of a legally protected interest which is (a) concrete and particularized, ... and (b) `actual or imminent, not ``conjectural'' or ``hypothetical,'' ' ... Second, there must be a causal connection between the injury and the conduct complained of---the injury has to be `fairly ... trace[able] to the challenged action of the defendant, and not ... th[e] result [of] the independent action of some third party not before the court.' ... Third, it must be `likely,' as opposed to merely `speculative,' that the injury will be `redressed by a favorable decision.' ''
\end{quote}
\noindent --- Justice Scalia, for the Court, 504 U.S. at 560--61.

\begin{quote}
``Such `some day' intentions---without any description of concrete plans, or indeed even any specification of when the some day will be---do not support a finding of the `actual or imminent' injury that our cases require.''
\end{quote}
\noindent --- Justice Scalia, for the Court, 504 U.S. at 564.

\begin{quote}
``In my view, Congress has the power to define injuries and articulate chains of causation that will give rise to a case or controversy where none existed before, and I do not read the Court's opinion to suggest a contrary view.''
\end{quote}
\noindent --- Justice Kennedy, concurring in part and concurring in the judgment, 504 U.S. at 580.

\end{document}
